% =====================================================================
% Chapter 6 -- Models, Training, and Hyperparameter Optimisation
% =====================================================================
\documentclass[../preamble.tex]{subfiles}
\begin{document}

\chapter{Models, Training, and Hyperparameter Optimisation}
\label{ch:training}

\section{Architecture roster}
\label{sec:train-roster}

The thesis uses nine ImageNet-pretrained backbones, all accessed through the
\texttt{timm} library and all standardised to a $224 \times 224$ input. The
roster spans residual, dense, mobile, convolutional-modern, and transformer
families. Their parameter counts range from 5.3M (EfficientNet-B0) to 86.6M
(ViT-Base/16).

\begin{table}[htbp]
  \centering
  \small
  \caption{Architecture roster. The \texttt{freeze\_fraction} column records the value used in the executed top-3 fine-tune; Optuna selected \texttt{freeze\_fraction=0} for all three members (see~\cref{tab:hpo-best}).}
  \label{tab:model-roster}
  \begin{tabular}{llrrl}
    \toprule
    \textbf{Name} & \textbf{timm key} & \textbf{Params (M)} & \textbf{Trainable (M)} & \textbf{Freeze frac.} \\
    \midrule
    ResNet-50 & \texttt{resnet50} & 25.6 & 25.6 & 0 \\
    ResNet-101 & \texttt{resnet101} & 44.5 & 44.5 & 0 \\
    DenseNet-121 & \texttt{densenet121} & 8.1 & 8.1 & 0 \\
    DenseNet-169 & \texttt{densenet169} & 14.1 & 14.1 & 0 \\
    EfficientNet-B0 & \texttt{tf\_efficientnet\_b0} & 5.3 & 5.3 & 0 \\
    ConvNeXt-Tiny & \texttt{convnext\_tiny} & 28.6 & 28.6 & 0 \\
    MobileNetV3-Large & \texttt{mobilenetv3\_large\_100} & 5.4 & 5.4 & 0 \\
    ViT-Base/16 & \texttt{vit\_base\_patch16\_224} & 86.6 & 86.6 & 0 \\
    Swin-Tiny & \texttt{swin\_tiny\_patch4\_window7\_224} & 28.3 & 28.3 & 0 \\
    \bottomrule
  \end{tabular}
\end{table}


\section{Classification head}
\label{sec:train-head}

Each backbone is followed by a SiLU-activated head with two hidden layers and
BatchNorm:
\begin{verbatim}
GlobalAvgPool -> Linear(in_features, 256) -> SiLU
             -> BatchNorm1d(256) -> Dropout(p)
             -> Linear(256, 128) -> SiLU
             -> BatchNorm1d(128)
             -> Linear(128, num_classes)
\end{verbatim}
The head reads \texttt{num\_classes} from the configuration rather than
hard-coding it. SiLU is used in place of ReLU because its smooth gradient
near zero is consistent with the EMA-style weight averaging that is
described in the methodology specification.

\section{Foundation-stage training}
\label{sec:train-foundation}

The foundation stage trains every backbone under both denoisers with default
hyperparameters, no HPO, and a single 80/10/10 split. The learning rate is
$1 \times 10^{-4}$ with a one-epoch warmup, weight decay $1 \times 10^{-3}$,
batch size 32, patience 5, and a maximum of 15 epochs. The internal AUC
plateaus near 1.0 across all eighteen configurations, which is why the
foundation stage alone is insufficient to characterise the pipeline.

\section{Fine-tune-stage training}
\label{sec:train-finetune}

The fine-tune stage starts from foundation checkpoints and adapts them to the
PCOSgen training pool. The learning rate is reduced to $3 \times 10^{-5}$
(one tenth of the foundation rate), weights-only state is reloaded into a
fresh optimiser, and training runs for up to 15 epochs with the same
patience of 5. The held-out PCOSgen test set is never used during fine-tune.

\begin{figure}[htbp]
  \centering
  \includegraphics[width=\linewidth]{figures/thesis_training_curves.png}
  \caption{Training and validation curves for the three ensemble members. The
           fine-tunes converge within a few epochs; the divergence in
           validation AUC by epoch 11 marks the best-epoch selection event.}
  \label{fig:thesis-training-curves}
\end{figure}

\section{Hyperparameter optimisation}
\label{sec:train-hpo}

Optuna was used for per-architecture HPO \cite{akiba2019}. Each backbone
started from its foundation checkpoint and ran 20 trials of 8 epochs each,
with MedianPruner cutting trials at epoch 4 if the trial was clearly
underperforming. The search space is:

\begin{itemize}
  \item learning rate: $[5 \times 10^{-6},\, 1 \times 10^{-4}]$, log-uniform
  \item weight decay: $[1 \times 10^{-5},\, 1 \times 10^{-2}]$, log-uniform
  \item dropout: $[0,\, 0.5]$
  \item rotation: $[0,\, 20]\,\text{degrees}$, integer
  \item label smoothing: $[0,\, 0.1]$
  \item freeze fraction: $\{0,\, 0.5\}$
  \item batch size: $\{16,\, 32,\, 64\}$
\end{itemize}

The full search-space inventory and the selected best parameters for the
three ensemble members are recorded in \cref{app:hpo-search-space}. The
table below summarises the best params that were selected:

\begin{table}[htbp]
  \centering
  \small
  \caption{Optuna best hyperparameters for the three ensemble members. All three selected \texttt{freeze\_fraction=0} (full backbone fine-tuning).}
  \label{tab:hpo-best}
  \begin{tabular}{lcccccc}
    \toprule
    \textbf{Architecture} & \textbf{LR} & \textbf{WD} & \textbf{Dropout} & \textbf{Rot.} & \textbf{LS} & \textbf{Batch} \\
    \midrule
    \bottomrule
  \end{tabular}
\end{table}


The key observation is that Optuna selected \texttt{freeze\_fraction=0} for
all three ensemble members, contradicting the v3 methodology default of
\texttt{freeze\_fraction=0.60}. The intended-versus-executed table in
\cref{app:intended-executed} documents this deviation explicitly. The
deviation is consistent with the observation that the external cohort
introduces new feature statistics that benefit from full backbone
adaptation.

\section{Training stability}
\label{sec:train-stability}

The trainer uses bf16 autocast and channels-last memory layout. The
divergence guard is implemented as an \texttt{optuna.TrialPruned} exception
triggered when a non-finite loss is encountered; this prevents a single
numerically bad batch from corrupting a sweep trial. Gradient clipping at
L2-norm 1.0 is applied per step. The methodology specification describes EMA
$\beta=0.999$, self-contained checkpoint files, and SWA over the last five
checkpoints. These features were implemented in the trainer scaffold but
were not used for the final reported results, as documented in
\cref{app:designed-not-executed}.

\section{Summary}
\label{sec:train-summary}

The training pipeline produces three high-AUC fine-tuned checkpoints,
each with a different inductive bias. The next chapter describes how they
are combined into a calibrated and uncertainty-aware ensemble.

\end{document}